% ======================================================================
%  Slow Loops and Fast Resonances: Hannay's Angle vs. Parametric Resonance
%  ----------------------------------------------------------------------
%  Standalone-compilable chapter. If you \input{} this into an existing
%  manuscript, keep only the material between
%      %%% BEGIN CHAPTER %%%   and   %%% END CHAPTER %%%
%  and make sure your master preamble loads: amsmath, amssymb, amsthm,
%  mathtools, booktabs. Everything else here is self-contained.
% ======================================================================
\documentclass[11pt]{book}

\usepackage[T1]{fontenc}
\usepackage{amsmath,amssymb,amsthm,mathtools}
\usepackage{booktabs}
\usepackage[margin=1in]{geometry}
\usepackage{enumitem}
\usepackage{hyperref}
\hypersetup{colorlinks=true,linkcolor=blue,citecolor=blue,urlcolor=blue}

\newcommand{\dd}{\mathrm{d}}
\newcommand{\ii}{\mathrm{i}}
\newcommand{\avg}[1]{\left\langle #1 \right\rangle}
\newcommand{\ord}[1]{\mathcal{O}\!\left(#1\right)}

\begin{document}

%%% BEGIN CHAPTER %%%
\chapter[Hannay's Angle versus Parametric Resonance]{Slow Loops and Fast Resonances: Hannay's Angle versus Parametric Resonance}
\label{ch:hannay-parametric}

\emph{Synopsis. Two classical results -- the exponential instability of a parametrically driven
oscillator, and the quiet geometric phase shift known as Hannay's angle -- look like they belong
to unrelated corners of mechanics. Written in action--angle variables, they turn out to be two
limiting regimes of a single perturbation series: the same Fourier harmonic of the same
perturbing term, treated in the adiabatic limit in one case and the resonant limit in the other.
This chapter derives both results from that common starting point, works through the standard
examples (the Mathieu oscillator, the Foucault pendulum), and is explicit about exactly where the
family resemblance ends.}

\section{Introduction}

Parametric resonance is the familiar fact that a periodically breathing oscillator -- a swing
whose rider stands up and sits down at the right moments, a pendulum whose length is pumped in
and out, an electrical circuit whose capacitance is modulated -- can absorb energy without ever
being pushed directly, provided the modulation frequency $\Omega$ sits near twice its natural
frequency $\omega_0$. Sufficiently close to that resonance, the amplitude grows exponentially;
sufficiently far from it, nothing much happens.

Hannay's angle is a far quieter effect. Take an integrable system whose Hamiltonian depends on
some external parameters, and carry those parameters slowly around a closed loop back to their
starting values. The action variable, being an adiabatic invariant, returns to where it started.
The angle variable does \emph{not}: on top of the ordinary dynamical phase $\int \omega\,\dd t$ it
picks up an extra shift that depends only on the geometry of the loop in parameter space, not on
how quickly the loop was traversed. It is the classical shadow of Berry's geometric phase in
quantum mechanics.

One of these is a runaway instability; the other is a bookkeeping correction so gentle that it is
easy to miss entirely. And yet, as the reader may already suspect from the way both are usually
introduced, something about the system's Hamiltonian has to change \emph{periodically} in both
cases for anything to happen at all. In the textbook parametric oscillator that periodic change is
imposed from the outside, at a driving frequency $\Omega$ set by an external agent. In the Hannay
problem the periodicity can equally well be imposed from outside -- an external knob traversing a
loop -- or arise \emph{internally}, generated by another, slower degree of freedom that belongs to
the very same autonomous system. This chapter shows that both phenomena are governed by the same
Hamiltonian perturbation series in action--angle variables, differing only in which harmonic of
that series is allowed to become secular, and that the "external vs.\ internal" distinction the
reader may have in mind maps onto a precise, and rather pretty, topological fact about the
parameter space itself.

\section{A Common Language: Action--Angle Variables with a Moving Parameter}
\label{sec:common}

\subsection{Action and angle, briefly}

For an integrable one-degree-of-freedom system with Hamiltonian $H_0(q,p;\lambda)$ depending on a
fixed parameter (or set of parameters) $\lambda$, the motion at fixed energy is periodic (or
quasiperiodic) in $q,p$. The action
\begin{equation}
I(E,\lambda) = \frac{1}{2\pi}\oint p \,\dd q
\end{equation}
and its canonically conjugate angle $\theta$ (defined mod $2\pi$) turn the dynamics into the
trivial form
\begin{equation}
H_0 = H_0(I;\lambda), \qquad \dot I = 0, \qquad \dot\theta = \omega(I;\lambda) \equiv
\frac{\partial H_0}{\partial I}.
\end{equation}
The transformation $(q,p)\leftrightarrow(I,\theta)$, for fixed $\lambda$, is generated by a
function $S(q,I;\lambda)$ with $p = \partial S/\partial q$, $\theta = \partial S/\partial I$.

\subsection{Letting $\lambda$ move}

Now let $\lambda = \lambda(t)$. The generating function $S$ becomes a time-dependent canonical
transformation, and the standard rule for such transformations -- new Hamiltonian equals old
Hamiltonian re-expressed in the new variables, plus $\partial S/\partial t$ -- gives
\begin{equation}
K(I,\theta,t) = H_0\big(I;\lambda(t)\big) + \dot\lambda \cdot \frac{\partial S}{\partial\lambda},
\label{eq:newham}
\end{equation}
since $S$ depends on $t$ only through $\lambda(t)$. The equations of motion generated by $K$ are
\begin{align}
\dot I &= -\frac{\partial K}{\partial\theta} = -\dot\lambda\cdot
\frac{\partial^2 S}{\partial\theta\,\partial\lambda}, \label{eq:Idot}\\[2mm]
\dot\theta &= \frac{\partial K}{\partial I} = \omega(I;\lambda(t)) + \dot\lambda\cdot
\frac{\partial^2 S}{\partial I\,\partial\lambda}. \label{eq:thetadot}
\end{align}
Everything in this chapter is a consequence of these two innocuous-looking equations, together
with the fact that $\partial S/\partial\lambda$, at fixed $I,\lambda$, is a $2\pi$-periodic
function of $\theta$ and therefore has a Fourier expansion
\begin{equation}
\frac{\partial S}{\partial\lambda}(I,\theta,\lambda) = g(I,\lambda) + \sum_{n\neq 0}
s_n(I,\lambda)\, e^{\ii n\theta}, \qquad g \equiv \avg{\partial S/\partial\lambda}_\theta,
\label{eq:fourier}
\end{equation}
where $\avg{\cdot}_\theta$ denotes the average over one period of $\theta$ at fixed $I,\lambda$.

\subsection{Two regimes of one perturbation series}
\label{sec:tworegimes}

Everything now hinges on how fast $\lambda(t)$ moves compared to the internal frequency
$\omega(I)$.

\begin{itemize}[leftmargin=1.4em]
\item If $\lambda(t)$ changes \textbf{slowly} -- $\dot\lambda/\omega \sim \varepsilon \ll 1$, with
no requirement that it be periodic in time at all, only that it eventually return to its starting
value, tracing a closed loop $C$ over some long time $T$ -- then in Eq.~\eqref{eq:thetadot} the
oscillatory harmonics $s_n\,e^{\ii n\theta}$ ($n\neq 0$) average to zero over one fast period and
contribute nothing secular; only the $n=0$ term $g(I,\lambda)$ survives as a steady drift. This is
the regime of \textbf{Hannay's angle} (\S\ref{sec:hannay}).

\item If $\lambda(t)$ oscillates \textbf{periodically} at a frequency $\Omega$ that is comparable
to $\omega(I)$ -- specifically, close to a low-order rational multiple, $\Omega \approx
\frac{n}{m}\omega(I)$ -- then it is instead one particular \emph{non-zero} harmonic, the one whose
combined phase $n\theta - m\Omega t$ varies slowly, that stops averaging away and becomes secular.
This is the regime of \textbf{parametric resonance} (\S\ref{sec:paramres}).
\end{itemize}

In both cases the action $I$ that would otherwise be a rigorous adiabatic invariant is disturbed
only by the surviving secular harmonic: a harmless, rate-independent phase shift in the first case,
an exponentially growing (or oscillating) amplitude in the second. The two phenomena are, in this
precise sense, the diagonal and an off-diagonal term of the very same object.

\section{Hannay's Angle}
\label{sec:hannay}

\subsection{Adiabatic invariance of $I$}

Averaging Eq.~\eqref{eq:Idot} over one fast period at fixed $I,\lambda$ kills every term in
Eq.~\eqref{eq:fourier} except none at all, since $\partial S/\partial\theta\,\partial\lambda$ has
zero $\theta$-average by construction (it is a pure derivative of a periodic function). Hence
$\avg{\dot I} = 0$ to this order: $I$ performs a small, bounded, purely oscillatory wobble as
$\lambda$ moves, with no net drift. This is the classical adiabatic theorem. It can be sharpened
considerably: for $\lambda(t)$ smooth (real-analytic), the actual drift in $I$ after one adiabatic
excursion is not merely small in $\varepsilon$ but exponentially small, $\sim
e^{-c/\varepsilon}$, a result due to Neishtadt (1975). For the purposes of this chapter the
leading-order statement -- $I$ returns to its initial value once $\lambda$ completes its loop --
is all we need.

\subsection{The secular drift of $\theta$: Hannay's formula}

The action being (adiabatically) conserved, attention shifts to the angle. Equation
\eqref{eq:thetadot} says that $\theta$ accumulates the ordinary dynamical phase $\omega(I;\lambda(t))$
plus a correction $\dot\lambda\cdot\partial^2S/\partial I\partial\lambda$. Using
Eq.~\eqref{eq:fourier}, the oscillatory part of this correction averages away over the loop just
as it did for $I$, but the $n=0$ piece does not: it is $\dot\lambda \cdot \partial g/\partial I$,
and integrating it over the full excursion $C$ gives a net, non-dynamical contribution to the
angle,
\begin{equation}
\boxed{\ \Delta\theta_{\text{Hannay}} = \oint_C \dd\lambda \cdot \frac{\partial g}{\partial I}
= -\oint_C \dd\lambda \cdot \avg{\frac{\partial\theta}{\partial\lambda}}_{I} \ }
\label{eq:hannay}
\end{equation}
(the second form, in which $\avg{\partial\theta/\partial\lambda}_I$ is the angle-averaged rate at
which the torus itself rotates under a change of $\lambda$ at fixed $I$, is the one Hannay wrote
down in his original 1985 paper and is the standard citation). This is \emph{Hannay's angle}: a
correction to the phase of the motion that depends only on the closed path $C$ traced in parameter
space, not on the rate $\dot\lambda(t)$ at which it is traversed -- provided only that the traverse
is slow enough for the averaging argument above to apply.

\subsection{The Berry--Hannay correspondence}

The resemblance to Berry's quantum geometric phase,
\begin{equation}
\gamma_n(C) = \ii\oint_C \avg{n(\lambda)|\nabla_\lambda|n(\lambda)}\cdot\dd\lambda,
\end{equation}
is not a loose analogy but a precise classical limit: Hannay's angle is what Berry's phase becomes
for a classical action-angle torus in place of a quantum eigenstate, with
$\avg{\partial\theta/\partial\lambda}_I$ playing the role of the Berry connection
$\avg{n|\nabla_\lambda|n}$. When $\lambda$ has two or more components, Stokes' theorem lets one
rewrite the loop integral in Eq.~\eqref{eq:hannay} as the flux of a curvature two-form through any
surface $\Sigma$ bounded by $C$,
\begin{equation}
\Delta\theta_{\text{Hannay}} = -\int_\Sigma F, \qquad F_{ab} = \frac{\partial A_b}{\partial
\lambda^a} - \frac{\partial A_a}{\partial\lambda^b}, \qquad A_a \equiv
\avg{\frac{\partial\theta}{\partial\lambda^a}}_I,
\label{eq:curvature}
\end{equation}
exactly as the Berry phase is the flux of the Berry curvature. As with Berry's phase, the apparent
paradox -- $A$ looks locally like the gradient of $g(I,\lambda)$, so how can its curl be non-zero?
-- is resolved the same way in both cases: $g$ (equivalently, the generating function $S$) can be
chosen smooth and single-valued only \emph{locally} in parameter space, not necessarily globally
over a whole neighbourhood of the loop, and it is exactly this obstruction that both formulas are
measuring. We will not need the fine print of that argument here; see Hannay (1985) and Berry
(1985) for the careful treatment.

\subsection{When is the geometric phase actually non-zero?}
\label{sec:onedim}

It is tempting to think any closed excursion of a parameter produces some geometric phase. It does
not, and the condition under which it fails is worth stating precisely, because it is exactly what
separates the "external, single-knob" version of these problems from the "internal" one the reader
may have had in mind.

\paragraph{A single parameter on an interval: always zero.} Suppose $\lambda$ is a single real
number -- a stiffness, a length, a frequency -- ranging over some interval, as in the ordinary
textbook parametric oscillator $\omega^2(t) = \omega_0^2\big(1+h\cos\Omega t\big)$. However
elaborate the excursion $\lambda(t)$, as long as $A(\lambda) \equiv
\avg{\partial\theta/\partial\lambda}_I$ is a genuine single-valued function of $\lambda$ alone, it
possesses an antiderivative $G(\lambda) = \int^\lambda A(\lambda')\,\dd\lambda'$, and
\begin{equation}
\oint_C A(\lambda)\,\dd\lambda = G\big(\lambda(T)\big) - G\big(\lambda(0)\big) = 0
\end{equation}
identically, because $\lambda(T)=\lambda(0)$. \emph{A single real parameter that goes out and comes
back produces no adiabatic geometric phase at all, to this order}: an interval is simply connected,
and the loop integral of a genuine gradient around any closed path in a simply connected domain
vanishes by the fundamental theorem of calculus. This is why the "obvious" analogue of Hannay's
angle for the everyday Mathieu oscillator -- slowly breathing the stiffness in and out -- simply is
not there. Getting a non-zero effect out of an externally imposed knob requires either of the two
mechanisms below.

\paragraph{Genuine curvature: two or more parameters.} If $\lambda$ has (at least) two independent
components, the loop $C$ can enclose a genuinely two-dimensional patch of parameter space on which
the curvature $F_{ab}$ of Eq.~\eqref{eq:curvature} need not vanish, and $\Delta\theta_{\text{Hannay}}$
is the flux of that curvature through the enclosed area. This is the Foucault-pendulum mechanism
worked out in \S\ref{sec:foucault}.

\paragraph{Topology instead of curvature: a circular parameter.} There is a second, less obvious
way to get a non-zero answer even from a single real parameter, if that parameter is not an
ordinary real number but an \emph{angle} $\varphi$ (defined mod $2\pi$), and the excursion winds
once around it rather than going out and back. The argument above breaks down precisely because a
circle is not simply connected: $G(\varphi) = \int^\varphi A(\varphi')\,\dd\varphi'$ need not close
up after one winding, and
\begin{equation}
\oint_0^{2\pi} A(\varphi)\,\dd\varphi = 2\pi\avg{A}_\varphi
\end{equation}
is generically non-zero. This is not a curvature effect at all -- there need be no local "curl" of
anything -- it is a pure holonomy around a non-contractible loop, the classical-mechanics cousin of
an Aharonov--Bohm phase picked up circling a solenoid rather than of a magnetic flux passing
through a cap. As we will see in \S\ref{sec:internal}, this is exactly the mechanism available to
an \emph{internally} generated loop, because the natural "parameter" supplied by another degree of
freedom of the same autonomous system is precisely an angle variable, and angle variables are
circles by construction.

\subsection{Worked example: the Foucault pendulum as parallel transport}
\label{sec:foucault}

The cleanest classical illustration of the curvature mechanism -- and the one Hannay and Berry both
used in their original papers -- is the Foucault pendulum. Idealize it as a two-dimensional
harmonic oscillator (small oscillations of a pendulum bob about its equilibrium), and use as the
"slow parameter" not a stiffness but a \emph{direction}: the local vertical at the pendulum's
location, viewed as a point on the celestial sphere $S^2$ as seen from an inertial frame fixed to
the distant stars. Sitting at a fixed point of latitude $\Lambda$ on the rotating Earth, this
direction traces out exactly one circle of latitude $\Lambda$ on that sphere every sidereal day, a
closed loop $C$ in a genuinely two-dimensional parameter space (the sphere). Because the sphere has
constant Gaussian curvature $K=1$ (in units where its radius is $1$), the flux of curvature through
the polar cap bounded by $C$ is simply the enclosed solid angle,
\begin{equation}
\Delta\theta_{\text{Hannay}} = -\Omega_{\text{solid}}(C) = -2\pi\big(1-\sin\Lambda\big),
\end{equation}
which, taken mod $2\pi$, is precisely $2\pi\sin\Lambda$ -- Foucault's classical result that the
plane of swing precesses, relative to the ground, by $360^\circ\sin\Lambda$ every sidereal day, a
full revolution at the poles and none at the equator. The ratio of the two relevant frequencies
here, the Earth's rotation rate against a pendulum's swing rate, is of order $10^{-5}$, comfortably
deep in the adiabatic regime that Eq.~\eqref{eq:hannay} assumes.

It is worth being candid about what this derivation does and does not add. The Foucault effect at
fixed latitude is also, and more elementarily, an \emph{exact} result: in the frame corotating with
the Earth the pendulum's Hamiltonian is time-independent, its two circular normal modes split by
the Coriolis coupling into frequencies $\omega_0 \pm \Omega_{\oplus}\sin\Lambda$, and the beat
between them reproduces the same precession with no adiabatic approximation needed at all. That the
"exact" and the "adiabatic/geometric" calculations agree exactly, not just to leading order in the
small ratio $\Omega_\oplus/\omega_0$, is itself a reflection of a fact used again in
\S\ref{sec:isochronism}: the harmonic oscillator's frequency does not depend on its amplitude, so
several approximations that are only asymptotic for a generic system become exact for this
particular one. The value of the Hannay-angle derivation is not that it is more accurate here --
it isn't -- but that it exposes the \emph{geometric} content of the result (a solid angle, a
holonomy) that the normal-mode calculation buries inside a beat frequency.

\subsection{An ``internal'' Hannay angle: autonomous slow--fast systems}
\label{sec:internal}

Nothing in the derivation of \S\ref{sec:common}--\S\ref{sec:hannay} required $\lambda(t)$ to be
imposed from outside. Suppose instead the full system is autonomous, with one integrable "fast"
degree of freedom $(I,\theta)$ and one "slow" degree of freedom $(J,\varphi)$, coupled so that the
fast subsystem's Hamiltonian depends on the slow \emph{angle} $\varphi$ rather than on time
directly,
\begin{equation}
H = H_{\text{slow}}(J) + H_{\text{fast}}(I,\theta;\varphi).
\end{equation}
To leading order in the separation of time scales -- the classical, Hamiltonian analogue of the
Born--Oppenheimer approximation -- the slow subsystem simply rotates at its own frequency,
$\dot\varphi \approx \Omega(J)$, playing exactly the role that the externally imposed $\lambda(t)$
played before, but now generated entirely from within the system itself: no external agent, no
clock, just the system's own slower motion acting as a moving parameter for its own faster motion.
This is precisely the "internal" periodicity referred to in the introduction.

And it is precisely the situation of the second bullet in \S\ref{sec:onedim}: $\varphi$ is not a
knob ranging over an interval, it is an angle, genuinely circular, and it winds around exactly once
per slow period. A single internal degree of freedom can therefore generate a non-zero Hannay angle
in the fast motion all by itself, with no need for a second, externally supplied parameter -- the
topology of the angle variable does the job that curvature had to do in the Foucault-pendulum case.
The reaction of this effect back on the slow variables was worked out in generality by Berry and
Robbins (1993) under the name \emph{geometric magnetism}: the fast subsystem exerts on the slow one
an effective force with exactly the structure of a Lorentz force, sourced by the same curvature (or
here, holonomy) that produces the Hannay angle in the first place. This is the classical shadow of
the geometric ("molecular Aharonov--Bohm") forces that appear in the Born--Oppenheimer treatment of
molecules, and it is a genuinely active area connecting classical mechanics, molecular physics, and
the geometric-phase literature; we leave the worked-out details to that literature and record here
only the structural point: \emph{internal} periodicity buys a mechanism for a non-zero adiabatic
geometric phase -- winding around a circle -- that is simply unavailable to a single
\emph{externally} imposed real-valued knob.

\section{Parametric Resonance}
\label{sec:paramres}

\subsection{Hill's equation and Floquet exponents}

The generic linear parametric oscillator is Hill's equation, $\ddot x + \omega^2(t)x = 0$ with
$\omega^2(t+2\pi/\Omega) = \omega^2(t)$ periodic. Floquet's theorem guarantees that the two
independent solutions take the form $x(t) = e^{\mu t}p(t)$ with $p$ periodic and $\mu$ (the Floquet
exponent) either purely imaginary -- bounded, quasiperiodic motion -- or possessing a non-zero real
part, in which case one solution grows without bound. The stability chart of this equation, plotted
against the modulation depth and the ratio $\Omega/\omega_0$, is a sequence of unstable "tongues"
emanating from the resonances $\Omega = 2\omega_0/n$, $n=1,2,3,\dots$, of which the $n=1$
("principal") resonance $\Omega\approx 2\omega_0$ is by far the widest and the one realized by
pumping a swing or breathing a pendulum's length.

\subsection{The same perturbation series, a different harmonic}

Consider the simplest concrete case, $\omega^2(t) = \omega_0^2\big(1+h\cos\Omega t\big)$, i.e.\ the
Hamiltonian $H = \tfrac12 p^2 + \tfrac12\omega_0^2\big(1+h\cos\Omega t\big)q^2$. Writing this in the
unperturbed harmonic-oscillator action--angle variables $q=\sqrt{2I/\omega_0}\,\sin\theta$,
$p=\sqrt{2I\omega_0}\,\cos\theta$ gives
\begin{equation}
H = I\omega_0 + h\,\omega_0 I\cos(\Omega t)\sin^2\theta
= I\omega_0 + \frac{h\omega_0 I}{2}\cos\Omega t - \frac{h\omega_0 I}{4}\Big[\cos(2\theta-\Omega t)
+ \cos(2\theta+\Omega t)\Big].
\label{eq:paramham}
\end{equation}
Exactly as anticipated in \S\ref{sec:tworegimes}, the perturbation is a sum of harmonics
$e^{\ii(n\theta - m\Omega t)}$; here $n=\pm2$, $m=\pm1$ dominate (the $\theta$-independent term
merely wiggles $\theta$ without consequence). Ordinary canonical perturbation theory removes such a
harmonic by dividing its coefficient by $n\omega(I)-m\Omega$ -- but for the harmonic oscillator
$\omega(I)=\omega_0$ is independent of $I$, so this denominator, $2\omega_0-\Omega$, is either
comfortably non-zero for all $I$ at once, or vanishes for all $I$ at once. There is no room to be
"a little bit resonant": either $\Omega$ is far from $2\omega_0$ and Eq.~\eqref{eq:paramham} is
perfectly well approximated by an adiabatic-style average with no secular effect on $I$ at all
(consistent with \S\ref{sec:onedim}: a single real parameter oscillating on an interval, away from
resonance, produces nothing), or $\Omega\approx 2\omega_0$ and the $n=2,m=1$ harmonic, whose
combined phase $\psi \equiv 2\theta - \Omega t$ is now slowly varying rather than fast, refuses to
average away and dominates the dynamics instead.

\subsection{Worked example: the growth rate}
\label{sec:growthrate}

Near the principal resonance it is simplest to work directly with the equation of motion. Set
$\Omega = 2\omega_0 + \delta$ with $\delta$ small, and write the solution as a slowly modulated
carrier, $x(t) = A(t)e^{\ii\omega_0 t} + A^*(t)e^{-\ii\omega_0 t}$ with $A$ slowly varying (the
method of averaging / rotating-wave approximation). Substituting into $\ddot x +
\omega_0^2(1+h\cos\Omega t)x = 0$ and keeping only the terms that resonate with $e^{\ii\omega_0 t}$
gives, after collecting coefficients,
\begin{equation}
2\ii\omega_0\dot A + \frac{h\omega_0^2}{2}A^*e^{\ii\delta t} = 0.
\end{equation}
Substituting $A = a\,e^{\ii\delta t/2}$ to remove the explicit time dependence yields the linear
system
\begin{equation}
\dot a = -\frac{\ii\delta}{2}a + \frac{\ii h\omega_0}{4}a^*,
\end{equation}
whose growth rate, found by the standard $a,a^*\sim e^{st}$ ansatz, is
\begin{equation}
\boxed{\ s = \sqrt{\left(\frac{h\omega_0}{4}\right)^2 - \left(\frac{\delta}{2}\right)^2}\ }
\label{eq:growthrate}
\end{equation}
Instability -- real, growing $s$ -- occurs for $|\delta| < h\omega_0/2$, i.e.\ for $\Omega$ within
$h\omega_0/2$ of $2\omega_0$ on either side, with the fastest growth, $s_{\max} = h\omega_0/4$,
exactly at resonance. This is the standard result for the width and strength of the principal
parametric-resonance tongue (see Landau \& Lifshitz, \emph{Mechanics}, \S27).

\subsection{Why the harmonic oscillator overreacts}
\label{sec:isochronism}

The unbounded exponential growth in Eq.~\eqref{eq:growthrate} is a special feature of the harmonic
oscillator's \emph{isochronism} -- the fact that $\omega(I)$ does not depend on $I$ at all. For a
generic (anharmonic) oscillator, $\omega(I)$ varies with amplitude, so as the resonantly pumped
amplitude grows, the system detunes itself out of resonance: $2\omega(I)-\Omega$ drifts away from
zero, halting the growth. Reduced to the pair $(I,\psi)$ with $\psi = 2\theta-\Omega t$ the slow
resonant angle, the dynamics there is that of a pendulum,
\begin{equation}
K(I,\psi) \approx \tfrac12 M_{\text{eff}}(I-I_\ast)^2 + \varepsilon V\cos\psi,
\end{equation}
with stable (elliptic) and unstable (hyperbolic, separatrix) fixed points -- a bounded "resonance
island" in phase space rather than an unbounded runaway. Only in the degenerate, isochronous limit
$M_{\text{eff}}\to 0$ does the island's width diverge and the bounded libration become the familiar
unbounded exponential instability of the linear Mathieu equation. It is amusing, and not an
accident, that motion \emph{around} such a resonance island -- itself a slow, closed orbit in the
reduced $(I,\psi)$ phase plane -- can in principle support its own second-generation Hannay angle
if some yet slower parameter is in turn carried around a loop; we will not pursue that here, but it
is a reminder that "fast" and "slow" are relative labels, not absolute ones.

\section{Putting Them Side by Side}

\subsection{One sentence}

If the whole of this chapter had to be compressed into a single sentence, it would be this: both
effects come from the same Fourier decomposition of the same perturbing term in $H(I,\theta;\lambda(t))$,
and the only question that matters is whether it is the \textbf{zeroth} harmonic (the
$\theta$-average, always secular, however slow $\lambda$ moves) or a \textbf{non-zero} harmonic
(secular only when tuned into resonance with the drive) that survives averaging. Hannay's angle is
the physics of the diagonal term; parametric resonance is the physics of an off-diagonal term
caught at resonance.

\subsection{A table of correspondences}

\begin{center}
\renewcommand{\arraystretch}{1.3}
\begin{tabular}{@{}p{3.1cm}p{5.7cm}p{5.7cm}@{}}
\toprule
& \textbf{Hannay's angle} & \textbf{Parametric resonance} \\
\midrule
Regime & $\Omega \ll \omega(I)$ (adiabatic) & $\Omega \approx \tfrac{n}{m}\,\omega(I)$ (resonant) \\
Relevant harmonic & $n=0$ (the $\theta$-average) & $n\neq 0$, matched to the drive \\
Effect on the action $I$ & Adiabatically invariant; bounded wobble only & Secular: unbounded growth
(harmonic case) or bounded libration in a resonance island (anharmonic case) \\
Effect on the angle $\theta$ & Extra phase shift per loop, independent of the traversal rate &
Resonant combination $n\theta-m\Omega t$ locks or librates instead of rotating \\
Governing structure & Berry--Hannay connection / curvature on parameter space &
Floquet exponents; Mathieu stability chart; Arnold tongues \\
Minimum parameter content & Two parameters (curvature), \emph{or} one topologically circular
parameter (holonomy) & One real parameter suffices \\
Sensitivity to rate & None, in the adiabatic limit: purely geometric & Essential: growth rate scales
with modulation depth $h$ \\
Canonical example & Foucault pendulum / slowly rotating anisotropic oscillator & Pumped swing /
Mathieu oscillator / parametric amplifier \\
Quantum shadow of & Berry's geometric phase & Floquet quasi-energy resonances (e.g.\ parametric
down-conversion, squeezing) \\
\bottomrule
\end{tabular}
\end{center}

\subsection{A remark on the quantum mirror}

The last row of the table is worth a moment's pause, because it shows the analogy running in two
directions at once. Hannay's angle is the classical limit of Berry's adiabatic phase; a
periodically driven quantum oscillator, on the other hand, has its own exact Floquet theory, in
which quasi-energies replace ordinary energies and resonant coupling between Floquet sidebands
plays the role that $2\omega_0\approx\Omega$ plays classically -- this is precisely the mechanism
behind quantum parametric amplification and squeezed-light generation. Classical mechanics thus
sits between two quantum phenomena that are not usually mentioned in the same breath: adiabatic
transport on the one hand, resonant multi-photon-like coupling on the other. That both descend from
the same classical perturbation series, differing only in which harmonic is allowed to go secular,
is a fair summary of this whole chapter.

\section{Closing Remarks}

The lesson is not that Hannay's angle and parametric resonance are secretly the same phenomenon --
they manifestly are not, one being a bounded geometric bookkeeping term and the other an
instability that can shake a system to pieces. The lesson is that both are limiting cases of one
and the same object, a periodically or slowly modulated integrable Hamiltonian written in
action--angle variables, and that the sharp physical distinction between them is entirely a matter
of which regime of a single small-denominator perturbation series one happens to be in. This is, in
miniature, an instance of a point made elsewhere in this book about the geometric formulation of
mechanics: quantities like $\avg{\partial\theta/\partial\lambda}_I$ are not artifacts of a
particular coordinate calculation but live intrinsically on the space of parameters, exactly as a
connection lives on a bundle independent of any local trivialization -- and it is precisely because
Hannay's angle is such a genuinely geometric object that it survives, untouched, in the one regime
(adiabatic, off-resonance) where the far more dramatic but far less geometric phenomenon of
parametric resonance has nothing at all to say.

\bigskip
\noindent\textbf{Further reading.}
\begin{itemize}[leftmargin=1.4em,itemsep=2pt]
\item J.~H.~Hannay, ``Angle variable holonomy in adiabatic excursion of an integrable Hamiltonian,''
\emph{J.\ Phys.\ A} \textbf{18}, 221 (1985) -- the original derivation, with the bouncing-ball,
oscillator-loop, and Foucault-pendulum examples.
\item M.~V.~Berry, ``Classical adiabatic angles and quantal adiabatic phase,'' \emph{J.\ Phys.\ A}
\textbf{18}, 15 (1985) -- the companion paper establishing the correspondence with the quantum
geometric phase.
\item M.~V.~Berry and J.~M.~Robbins, ``Chaotic classical and half-classical adiabatic reactions:
geometric magnetism and deterministic friction,'' \emph{Proc.\ R.\ Soc.\ Lond.\ A} \textbf{442}, 659
(1993) -- the internal, autonomous slow--fast version and the resulting geometric force.
\item L.~D.~Landau and E.~M.~Lifshitz, \emph{Mechanics}, 3rd ed., \S27 -- the standard treatment of
parametric resonance and the Mathieu equation.
\item V.~I.~Arnold, \emph{Mathematical Methods of Classical Mechanics}, 2nd ed. -- action--angle
variables, adiabatic invariants, and resonance in the broader Hamiltonian context.
\end{itemize}

%%% END CHAPTER %%%
\end{document}
